\section{Integrating Factors \& Separable Equations}

\subsection{More Examples of Integrating Factors}

\begin{example}
    Solve the differential equation $\frac{dy}{dt} - y = e^{-t}$.
\end{example}
\begin{solution}
    Start by multiplying the equation by the integrating factor $\mu(t)$:
    \begin{align}
        \mu(t)\frac{dy}{dt} - \mu(t)y = \mu(t)e^{-t}
    \end{align}
    Using Remark 2.2.3. (remember this), we find the integrating factor. Let $p(t) = -1$: 
    \begin{align*}
        &\hspace{1.2cm} \mu(t) = \exp \int -dt \\
        &\iff \mu(t) = e^{-t + C} \\
        &\iff \mu(t) = Ae^{-t} 
     \end{align*}
    Equation $(4)$ becomes:
    \begin{align}
        e^{-t}\frac{dy}{dt} - e^{-t}y = e^{-2t}
    \end{align}
    We recognize the left side of equation $(5)$ as the derivative of the product $e^{-t}y$. Integrating both sides, we obtain:
    \begin{align*}
        &\hspace{1cm}\cancel{\int \frac{d}{dt}}\left[e^{-t}y\right]\cancel{dt} = \int e^{-2t}dt \\
        &\implies e^{-t}y = -\frac{e^-2t}{2} + C \\
        &\implies y = -\frac{e^{-t}}{2} + Ce^t \qedhere
    \end{align*}
\end{solution}
\begin{example}
    Solve the differential equation $\frac{dy}{dt} = y\sin(t) + 2te^{-\cos(t)}$.
\end{example}
\begin{solution}
    Let us first rearrange the equation to put all terms with $y$ on the left side, then multiply by the integrating factor $\mu(t)$:
    \begin{align}
        \mu(t)\frac{dy}{dt} - \mu(t)y\sin(t) = \mu(t)2te^{-\cos(t)}
    \end{align}
    We recall Remark 2.2.3 (again, remember this) to find $\mu(t)$. Let $p(t) = -\sin(t)$:
    \begin{align*}
        \mu(t) &= \exp \int -\sin(t) dt \\
        &= e^{\cos(t) + C} \\
        &= Ae^{\cos(t)}
    \end{align*}
    Plugging our newfound value of $\mu(t)$ into Equation $(6)$:
    \[ e^{\cos(t)}\frac{dy}{dt} - e^{\cos(t)}y\sin(t) = 2t \]
    We recognize the left side of the equation as the derivative of the product $e^{\cos(t)}y$. Integrating and isolating for $y$ yields:
    \begin{align*}
        &\hspace{1.2cm} \cancel{\int \frac{d}{dt}}\left[e^{\cos(t)}y\right] \cancel{dt} = \int 2t dt\\
        &\iff e^{\cos(t)}y = t^2 + C \\
        &\iff y = e^{-\cos(t)}(t^2 + C) \qedhere
    \end{align*}
\end{solution}

\subsection{Separable Equations}

We learned how to use integrating factors to solve differential equations of the form $\frac{dy}{dt} + p(x)y = g(x)$. We will now shift our focus to differential equations of the form $\frac{dy}{dt} = f(y)g(t)$. We say this class of differential equations are \bold{separable}, and we can solve them by putting like variables on the same sides of the equality, and then integrating both sides.

\begin{example}[Newton's Law of Cooling]
    Consider the following ODE:
    \[\frac{dB}{dt} = K(E - B)\]
    where $E, K \in \RR$, with the initial condition $B(0) = B_0$. To solve this equation, we begin by putting variables with $B$ on the left side, and variables with $t$ (there are none) on the right side:
    \[\frac{dB}{E - B} = K dt\]
    Integrating both sides yields:
    \begin{align*}
        &\hspace{1cm} \int \frac{dB}{E - B} = \int K dt \\
        &\implies -\ln|E - B| = Kt + C \\
        &\implies \ln|E - B| = -Kt - C \\
        &\implies E - B = \pm e^{-Kt - C} \\
        &\implies -B = Ae^{-Kt} - E \\
        &\implies B = E - Ae^{-Kt} 
    \end{align*}
    where $A = \pm e^{-C}$. We have that $B(0) = B_0$, so we can find the value of $A$ given the initial value:
    \begin{align*}
        &\hspace{1cm} B_0 = E - A \\
        &\implies A = E - B_0
    \end{align*}
    Hence the solution to the ODE is:
    \[B(t) = E - (E - B_0)e^{-Kt}\]
\end{example} 

\begin{example}
    Solve the differential equation $\frac{dy}{dt} = 6y^2t$, given the initial condition $y(1) = \frac{1}{3}$. Additionally, find the interval of valid solutions for $t$.
\end{example}
\begin{solution}
    We will solve this equation by separation. First, bring all terms with $y$ on the left side and all with $t$ on the right side:
    \[\frac{dy}{y^2} = 6t dt\]
    Integrating both sides of the equation yields:
    \begin{align*}
        &\hspace{1cm} \int \frac{dy}{y^2} = \int 6t dt \\
        & \implies \frac{-1}{y} = 3t^2 + C \\
        & \implies y = \frac{-1}{3t^2 + C}
    \end{align*}
    Now, find $C$ using the initial value $y(1) = \frac{1}{3}$:
    \begin{align*}
        &\hspace{1cm} \frac{1}{3} = \frac{-1}{3 + C} \\
        & \implies 3 = -3 - C \\
        & \implies C = -6
    \end{align*}
    So the solution to the equation is:
    \[y = \frac{-1}{3t^2 - 6}\]
    Now to find valid solutions for $t$, we need to find the restrictions on the solution we found above. We know that $3t^2 - 6 \neq 0 \implies 3t^2 \neq 6 \implies t \neq \pm\sqrt{2}$. So the interval of valid solutions for $t$ is:
    \[t = \{t \in \RR \colon t \neq \pm \sqrt{2}\}\]
    This implies that the solutions are discontinuous at points $t = -\sqrt{2}$ and $t = \sqrt{2}$.
\end{solution}

\begin{example}
    Solve the differential equation $\frac{dy}{dt} = \frac{ty^3}{1 + t^2}$ with the initial condition $y(0) = 1$.
\end{example}
\begin{solution}
    We will solve by separation. We first bring all terms with $y$ on the left side and all with $t$ on the right:
    \[\frac{dy}{y^3} = \frac{t}{1 + t^2}dt\]
    Integrating both sides yields:
    \begin{align*}
        &\hspace{1cm} \int \frac{dy}{y^3} = \int \frac{t}{1 + t^2}dt \\
        & \implies \frac{-1}{2y^2} = \frac{1}{2}\ln|1 + t^2| + C \\
        &\implies y^{-2} = -\ln|1 + t^2| - 2C \\
        &\implies y^2 = \frac{1}{-\ln|1 + t^2| - 2C} \\
        &\implies y = \pm\frac{1}{\sqrt{-\ln|1 + t^2| - 2C}}
    \end{align*}
    Using the intial condition $y(0) = 1$ we solve for $C$:
    \begin{align*}
        &\hspace{1cm} 1 = \pm \frac{1}{\sqrt{-2C}} \\
        &\implies \pm\sqrt{-2C} = 1 \\
        &\implies -2C = 1 \\
        &\implies C = \frac{-1}{2}
    \end{align*}
    So the solution to the equation is:
    \[y = \frac{1}{\sqrt{1 - \ln|1 + t^2|}} \qedhere\] 
\end{solution}

\begin{example}
    Solve the differential equation $\frac{dy}{dt} = e^{y - t}\sec(y)(1 + t^2)$ with the intiial condition $y(0) = 0$. 
\end{example}
\begin{solution}
    Again, we proceed by separating $y$ and $t$:
    \begin{align*}
        &\hspace{1.2cm }\frac{dy}{e^y\sec(y)} = e^{-t}(1 + t^2)dt \\
        &\iff e^{-y}\cos(y) dy =  e^{-t}(1 + t^2)dt
    \end{align*}
    Integrating both sides can be done by using integration by parts on both integrals, though for brevity I will omit those calculations:
    \begin{align*}
        &\hspace{1cm} \int e^{-y}\cos(y) dy =  \int e^{-t}(1 + t^2)dt \\
        &\implies \frac{e^{-y}(\sin(y) - \cos(y))}{2} = e^{-t}(t^2 + 2t + 3) + C  \\
        &\implies e^{-y}(\sin(y) - \cos(y)) = 2e^{-t}(t^2 + 2t + 3) + 2C
    \end{align*}
    This is unfortunately the best we can simplify the equation, but we can still solve for $C$ using the condition $y(0) = 0$:
    \begin{align*}
        &\hspace{1cm} 1 = 6 + 2C \\
        &\implies C = \frac{-5}{2}
    \end{align*}
    So an implicit solution to the differential equation is:
    \[e^{-y}(\sin(y) - \cos(y)) = 2e^{-t}(t^2 + 2t + 3) - 5\]
\end{solution}



